\scanpage{6}
ou comme d'habitude l'indice $n$ indique la réduction mod \deleted{la puissance}
$\underline{m}^{n+1}$ ($\underline{m}$ étant l'idéal maximal de $A$). Pour tout entier $m\quad0$, dénotant
par ${}_m u,{}_m v$ les homomorphismes ${}_m H\quad G$ induits par $u,v$, on aura donc
aussi $({}_m u)_n=({}_m v)_n$. Ceci étant vrai pour tout $n$, et ${}_m H$ étant \uline{fini}
sur $S$ en vertu de 2.2. , il s'ensuit qu'on aura ${}_m u={}_m v$. Ceci étant
vrai pour tout $m$, on aura $u=v$ en vertu de 4.7. a).

b) Cas général. Nous \deleted{\unsure{supposerons}} pouvons supposer $S$ \deleted{local
\unsure{de point fermé $s$.}} \inserted{affine, $S=\operatorname{Spec}(A)$.} Quitte à faire une extension fidèlement plate et quasi-compacte de la base, on peut supposer de plus $H$ diagonalisable, donc
de la forme $D_S(M)$, avec $M$ un groupe commutatif de type fini. Introduisons comme dans la démonstration de 4.6. la famille filtrante croissante des sous-$\underline{Z}$-algèbres \deleted{$A_i$} de type fini $A_i$ de $A$. Comme $G$ est de présentation finie sur $S$, il provient d'un \deleted{schéma} préschéma en groupes de
présentation finie $G_i$ sur un $S_i$. D'ailleurs, $H$ provient de $H_i=D_{S_i}(M)$.
De plus, pour $j$ grand, en vertu de EGA IV 8, les homomorphismes
\deleted{$u,v$} $u,v$ proviennent d'homomorphismes de $S_j$-groupes \deleted{$\unsure{H}$}$u_j,v_j:H_j\quad G_j$,
\deleted{Soit $s_j$ \unsure{le}} et $u_j$ sera central pour $j$ grand. D'ailleurs, si $s_j$ est
l'image de $s$ dans $S_j$, on a $u_{j\,s_j}=v_{j\,s_j}$. On peut alors appliquer a)
à la situation sur $S_j$, d'où aussitôt la conclusion annoncée.

\uline{Corollaire} 5.2. Soient $u:H\quad G$ un homomorphisme de $S$-préschémas en
groupes, avec $H$ de type multiplicatif et de type fini, et $H$ de présentation finie (ou seulement localement de \deleted{\unsure{pr}} type fini sur $S$, lorsque
$S$ est localement noethérien). Soit $s\quad S$ et supposons que $u_s$ soit l'homomorphisme unité. Alors il existe un voisinage ouvert $U$ de $s$ tel que
$u_U$ soit l'homomorphisme unité.

On prouve de la même façon les résultats suivants (qui contiennent
d'ailleurs les précédents, comme on voit en introduisant des \deleted{\unsure{sous-gr}}
graphes des homomorphismes en groupes précédents):

\margin{5}

